% ------------------------------------------------------------------------------------------
\section*{Introduction}
Brain metastases are monitored with serial gadolinium-enhanced MRI, and repeated gadolinium
exposure motivates synthesizing post-contrast images from non-contrast ones. Single-exam
synthesis must infer enhancement from T1, T2 and FLAIR alone. In surveillance, however, every
patient has a prior post-contrast exam, highly correlated with the current one. We condition an
image-to-image Schr\"odinger bridge (I$^2$SB) on that prior exam and evaluate it under
inference-realistic conditions.

\section*{Methods}
\textbf{Image-to-Image Schr\"odinger Bridge:} I$^2$SB seeks to take some image $\bm{x}_1$,
sampled from some data distribution, and estimate some related $\bm{x}_0$ coming from some
alternative target distribution. Generating a sample $\bm{x}_t$ from the bridge involves
generating noisy convex combinations of $\bm{x}_0$ and $\bm{x}_1$, where the additive noise
power is time varying for $t \in [0, 1]$. A network is trained to, given bridge sample
$\bm{x}_t$, predict the bridge endpoint $\bm{x}_0$. The algorithm itself iteratively predicts
$\bm{x}_0$ and then resamples a new bridge point $\bm{x}_t$ while annealing the time index to $0$
(Figure~\ref{fig:i2sb}). In our case, we bridge from the current pre-contrast image to the
desired post-contrast image. The original I$^2$SB paper uses a UNet with an appended
time-dependent channel. Since we have access to other contrasts of the same anatomy, we
additionally append [T1, T2, FLAIR] as input channels, so the model need only learn the
enhancement pattern from T1 to CT1.\\

\textbf{Longitudinal Conditioning:} Where enhancement is not evident in the pre-contrast images,
the patient's prior contrast-enhanced exam can supply it. The bridge endpoints are unchanged
(current T1 to current CT1), and the prior exam enters only as conditioning: its T1 and CT1 at a
5-slice window are appended as ten further input channels. The neighboring slices robustify the
network to through-plane misregistration.\\

\textbf{Experiments:} We trained on NYUMets, a longitudinal dataset of brain-metastasis patients
with FLAIR, T1, CT1 and T2 at each visit, split into a training set of 2318 exams and a holdout
set of 1007 exams. Within each exam the contrasts are already co-registered and skull-stripped.
We train two regressors over the same training set while only giving one access to prior-exam
data. At inference, the prior exam is registered to the current one by a 3-D affine fit learned
only from pre-contrast images, without relying on the CT1 being predicted. PSNR, SSIM and NRMSE
were computed within the brain mask and averaged per exam, and paired Wilcoxon signed-rank tests
were run over the holdout set.

\section*{Results}
Across 1007 test exams, conditioning on the prior exam improved every metric
(Figure~\ref{fig:bars}): PSNR rose from $24.79 \pm 0.09$ to $25.73 \pm 0.10$\,dB,
SSIM from $87.79 \pm 0.21\%$ to $90.83 \pm 0.19\%$, and NRMSE fell from
$13.07 \pm 0.19\%$ to $11.97 \pm 0.19\%$ (mean $\pm$ s.e.\ over exams; all
$p < 10^{-3}$, paired Wilcoxon signed-rank test). The improvement was consistent
rather than driven by a few cases: the longitudinal model was better in 97\% of exams
(median PSNR gain $+0.84$\,dB). Appending input channels to the UNet comes at little
computational cost, so both methods are effectively computationally matched.

We can see strong visual improvement in Figure~\ref{fig:cases}. In the first row, the
single-study I$^2$SB result overpredicts the enhancement details whereas longitudinal I$^2$SB
better preserves anatomical detail. The second row gives an example where a lesion appears in
the CT1 image and not in the current T1. In this case, the single-study I$^2$SB result correctly
predicts the appearance of the lesion but the longitudinal I$^2$SB result more accurately
recovers its shape. Additionally, the prior CT1 guide seems to contain enhancement details near
the brain stem that do not appear in the ground truth; longitudinal I$^2$SB is robust to them.
In the third row, single-study I$^2$SB hallucinates a lesion absent from the acquired CT1.

Figure~\ref{fig:poor} shows further cases where interval change or misregistration makes the
prior a poor guide. Figure~\ref{fig:sim} quantifies this by scoring the registered prior CT1
against the current CT1: the gain grows from $\approx+0.3$\,dB for the least similar priors to
$+2.5$\,dB for the most similar, and stays positive on average in every bin.

\section*{Discussion}
The prior exam is routinely available in surveillance yet unused by contrast-synthesis methods.
As conditioning, it lets the network transfer known enhancement rather than infer it, and a
poor prior reduces the benefit without degrading synthesis on average. Limitations: image-level
metrics only, and new lesions, which the prior cannot inform, need lesion-level evaluation.

\section*{Conclusion}
Conditioning I$^2$SB on a patient's prior exam, registered without the scan being predicted,
improves gadolinium-free synthesis of the follow-up post-contrast T1.

% ------------------------------------------------------------------------------------------
% Figure labels used above: fig:i2sb (Fig 1), fig:cases (Fig 2), fig:bars (Fig 3),
% fig:poor (Fig 4), fig:sim (Fig 5).
